\documentclass[a4paper]{article}

\usepackage{graphicx}
\usepackage{amsmath,amssymb}
\usepackage{minted}
\usepackage{multirow}
\usepackage{float}
\usepackage{todonotes}
\usepackage{forest}
\usepackage{xstring}
\usepackage{xcolor}
\usepackage[most]{tcolorbox}
\usepackage{xparse}
\usepackage{natbib}

\usetikzlibrary{graphs,graphdrawing}
\usegdlibrary{layered}

% for external graphviz
\input{/home/donkarlo/Dropbox/repo/nd_ai_project/src/nd_ai/knoweldge_representation/language/formal/computer/programming/tex/latex/code_clips/external_graphviz.tex}

% for tags
\input{/home/donkarlo/Dropbox/repo/nd_ai_project/src/nd_ai/knoweldge_representation/language/formal/computer/programming/tex/latex/parts/control_sequence/kind/semtag/semtag.tex}

% for links
\input{/home/donkarlo/Dropbox/repo/nd_ai_project/src/nd_ai/knoweldge_representation/language/formal/computer/programming/tex/latex/code_clips/seehere.tex}

\title{self-organisation}
\date{}

\begin{document}
    \maketitle
    \cite{ashby-1947-principles-of-the-self-organizing-dynamic-system} presents one of the classic formulations of self-organization and argues that some dynamical systems can move toward order and stability without an external designer.
    His focus is mainly cybernetic and systemic: how a system, through its own internal dynamics, can settle into more organized states.
    This work lays a theoretical foundation for self-organization, rather than explaining life directly in the language of modern biology.
    
    Biological organization is the organization of complex biological structures and systems that define life using a reductionistic approach \seehere{https://en.wikipedia.org/wiki/Biological_organisation}.
    
    
    
    \cite{kauffman-1993-the-origins-of-order-self-organization-and-selection-in-evolution} argues that biological order does not arise from natural selection alone; self-organization also plays a fundamental role.
    He shows that complex networks can spontaneously generate patterns, structures, and stable constraints, and that these provide the basis on which evolution operates.
    In this view, life emerges from the interplay of selection and self-organization, not from only one of them.
    
    
    
    
    \bibliographystyle{plainnat}
    \bibliography{/home/donkarlo/Dropbox/repo/nd_shared_research/bibliography/bibliography.bib}
\end{document}